% =========================================================================
% 2. RELATED WORK   [~0.75 page]
% GOAL: show existing evaluations don't cover the gap -> justify a new metric.
%       Organise by what each bucket MEASURES and what it MISSES (not a laundry list).
% NOTE: full cited survey still pending (deep-research). Verify every cite in custom.bib.
% =========================================================================
\section{Related Work}
\label{sec:related}

% ---- ¶1 (Jailbreak / attack evals): StrongREJECT, HarmBench, JailbreakBench (we use),
%      AdvBench, BeaverTails (we use). What they measure: Potency (does a harmful
%      request succeed). StrongREJECT adds a quality-aware rubric. What they miss:
%      no safety-reasoning / monitorability axis; mostly black-box, no intervention
%      notion. This is the closest prior art to our Potency axis.
\subsection{Attack evaluation}
Several benchmarks have been proposed as a standard for LLM safety-evaluation.
StrongReject \citep{souly2024strongreject}, HarmBench \citep{mazeika2024harmbench} and JailbreakBench \citep{chao2024jailbreakbench} (which we use) all provide small datasets of queries curated to be especially harmful, while BeaverTails \citep{ji2023beavertails} (which we use) provides a comprehensive set of QA-pairs with potentially harmful queries.

HarmBench \citep{mazeika2024harmbench} and JailbreakBench \citep{chao2024jailbreakbench} report ASR using LLM-as-a-judge to classify if a response complies with the harmful query. StrongReject \citep{souly2024strongreject} also uses an LLM to classify compliance, but additionally asks to score if the response is specific and convincing and reports a metric composed of the 3. 

While these benchmarks deal with model-safety, our metric evaluates the effect of intervention, comparing responses before and after.
In the same vein as StrongReject, our metric reflects harmfulness and quality of responses, but adds monitorability as a factor, reflecting if the potential harm is verbalized in reasoning.

% ---- ¶2 (Safety-classifier judges): Llama-Guard-3, ShieldGemma, our 30B/14B judges.
%      These are instrumentation (a better ASR), not a cross-method comparison
%      framework. Position our judges as instruments *inside* the framework, validated.
\subsection{Safety-classifier judges}
To compute an ASR that is more general than refusal or harmful keywords, we are inspired by
LlamaGuard-3 \citep{inan2023llamaguard} and ShieldGemma \citep{zeng2024shieldgemma} - LLMs specifically developed to classify potentially harmful in-/output in an effort to mitigate safety-issues. In addition, \citep{menke2025annotating} use the general-purpose Gemini 2.0-flash \textcolor{red}{cite} to classify safety-related reasoning.

Our proposed framework uses general-purpose LLMs that we evaluate for the tasks. The LLMs as judges are integral parts of the Potency and Safety Reasoning axes. In addition, we fine-tune an LLM to classify safety-reasoning traces.

% ---- ¶3 (White-box intervention methods --- what each REPORTS, and on what axis):
%      safety-head attribution/ablation (SHIPS/Sahara, refusal-rate); directional
%      ablation (Arditi et al., refusal + ad-hoc coherence check); representation /
%      activation steering; DSH harmfulness/refusal axes (Wu et al.); learned circuit
%      masks. Recurring gap: bespoke single-axis metrics, no shared quality or
%      reasoning axis, no cross-method protocol. THIS is the gap we fill.
\subsection{White-box intervention methods}
The effect of intervention methods on safety is commonly reported using ASR as the single measurement. \citep{zhao2025neuron} report ASR using generation of specific harmful responses as a success-criteria, following \citep{zou2023universal}.
\citep{zhou2024ships} reports ASR using the absence of refusal keywords as success-criteria. \citet{arditi2024refusal} also reports this, but in addition reports a score using an LLM to address potential incoherent responses. \citep{yu2026safeseek} reports ASR using  LlamaGuard3-8B to classify if an attack is successful.

We propose a standard protocol for measuring how ASR is affected by an intervention, reflected in the Potency-axis. Furthermore, our metric also encapsulates whether quality and reasoning are affected.

% ---- ¶4 (CoT-monitoring agenda): OpenAI 2024 / Baker et al. --- the reasoning trace
%      as a monitorable safety signal. Motivates our Safety-Reasoning axis; but it is
%      not itself an intervention-comparison metric.
\subsection{CoT-monitoring}
\citep{baker2025cotmonitoring} shows that an LLM can be used to monitor the CoT for signals of harmful behavior that may not be caught in the answer. In the same vein, \citep{korbak2025monitorability} mentions that the reasoning LLMs use CoT in the same way as activations, and argues that monitoring reasoning is a valuable additional safety layer.

With the Safety-Reasoning axis we address whether an intervention affects harmful signals that can be monitored in the CoT.

% ---- ¶5 (Composite / multi-metric precedents we borrow from): HELM (multi-metric
%      reporting is legitimate), StrongREJECT's quality-gated score, Arena/Elo ranking.
%      These inform *how* we combine axes defensibly (geometric mean + Pareto, §3.2).
\subsection{Composite metrics}
\citep{liang2023helm} argues that several characteristics are desirable for LLM responses and so LLMs should not be evaluated on a single axis, but rather report multiple metrics. \citep{souly2024strongreject} provides a composite metric of safety, combining quality-gated responses with the intent to comply with harmful responses.

Our metric consists of 3 axes that we find important for evaluating how effectively intervention methods locate safety mechanisms. We report the vector along the 3 axes with a Pareto front and additionally define how these axes can be combined into a single number

\TODO{Related Work is outline-only. Replace with a cited survey once deep-research is
run; do NOT cite specifics until custom.bib entries are verified.}
